\documentclass[10pt,a4paper]{amsart}
\usepackage[margin=19mm]{geometry}
\usepackage{amsmath,amssymb,mathtools}
\usepackage[scr=rsfs,cal=euler]{mathalfa}
\usepackage{xfrac}
\usepackage[unicode,hidelinks,bookmarks=false]{hyperref}
\newcommand{\ie}{i.e.\ }
\newcommand{\cf}{cf.\ }
\newcommand{\resp}{respectively\ }
\newcommand{\Subsection}{\subsection}
\providecommand{\nobreakdash}{\nobreak\hbox{-}}
\setlength{\emergencystretch}{2em}
\multlinegap0pt
\usepackage{bm}
\usepackage{bbm}
\usepackage{dsfont}
\usepackage{xspace}
\usepackage{cancel}
\usepackage{mathtools}
\usepackage{amsthm}
\makeatletter
\@ifundefined{theorem}{\newtheorem{theorem}{Theorem}}{}
\@ifundefined{lemma}{\newtheorem{lemma}[theorem]{Lemma}}{}
\@ifundefined{proposition}{\newtheorem{proposition}[theorem]{Proposition}}{}
\makeatother
\newcommand{\CC}{{\mathbb C}}
\newcommand{\R}{{\mathbb R}}
\newcommand{\ed}{\stackrel{d}{=}}
\title[On the pair correlation function of the Sine$\_\beta$ process]{On the pair correlation function of the Sine$\_\beta$ process}
\author{Yahui Qu, Benedek Valk\'o}
\date{Source: arXiv:2509.15446v1 (CC BY 4.0)}
\begin{document}
\maketitle

\noindent\emph{Source and licence.} On the pair correlation function of the Sine$_\beta$ process. Version arXiv:2509.15446v1 (CC BY 4.0), distributed under the Creative Commons Attribution 4.0 International licence, as recorded in the arXiv licence metadata for this exact version and shown on the abstract page https://arxiv.org/abs/2509.15446v1. Full rights notice: licenses/arxiv-2509.15446-CC-BY-4.0.txt.\par\smallskip

\noindent\textit{Editorial scope.} Counted unit: the Theorem whose source label is \texttt{thm:even\_beta} in the article (source0.tex lines 953--962, source label \texttt{thm:even\_beta}), delivered together with the article's own complete original proof of it (source0.tex lines 1002--1058, one \texttt{proof} environment of 57 lines). No mathematics is written, repaired or re-derived: statement and proof are the article's own text line for line. The proof is detached from the statement in the source: the article states the result at lines 953--962 and prints its proof at lines 1002--1058, and the proof environment's own header names the statement, so the association is the article's own. Required context retained uncounted: eq:alphaSDE (lines 450--453); eq:SDE\_initial (lines 455--457); prop:alphaSDE (lines 481--499); eq:QkSDE (lines 674--677); eq:q\_coeff\_1 (lines 943--945); lem:coeff\_s (lines 1061--1071). Every definition, display and label of those lines is retained unchanged. Omitted from this package and not redistributed: 1--449, 454--454, 458--480, 500--673, 678--942, 946--952, 963--1001, 1059--1060, 1072--1698. Nothing inside a retained excerpt is omitted or altered: every retained body is delivered exactly as the article's lines are written, so no omission receipt of any kind accompanies this package. Citations: the retained excerpts carry no citation command at all, so no bibliography entry of the article is transcribed and no reference is redistributed. Reference adaptation: none was needed and none was made; every reference inside the delivered unit resolves inside it, so no unresolved reference or citation is produced. Printed slips kept literal: no printed word, symbol, hypothesis or number of the counted statement or of its proof was altered. Layout and class: the article's own class is \texttt{article} with packages including palatino, amsmath, amsthm, amssymb, graphicx, enumerate, color, url, hyperref; the wrapper is the lane's pinned \texttt{amsart} editorial wrapper at 10pt on A4, which restates the theorem-like environments the retained text needs and adds the article's own macro definitions that the retained text needs (\texttt{\textbackslash CC}, \texttt{\textbackslash R}, \texttt{\textbackslash ed}), with extra wrapper packages (bm, bbm, dsfont, xspace, cancel, mathtools). The delivered numbering is the unit's own. Assets: no retained span contains a figure environment, a graphics-inclusion command, a PDF-page inclusion, a table or a TikZ picture; no image file is copied into this package, the package declares no asset dependency and no image file is redistributed with it. Licence: version arXiv:2509.15446v1 of this article is distributed under the Creative Commons Attribution 4.0 International licence, as recorded in the arXiv licence metadata for this exact version and shown on the abstract page https://arxiv.org/abs/2509.15446v1. A full rights notice is delivered at licenses/arxiv-2509.15446-CC-BY-4.0.txt. Characters: every retained excerpt is pure ASCII, so the delivered engine, XeLaTeX with its default Latin Modern text font, reports no missing character.
\begin{align}
    \label{eq:alphaSDE}
	d\alpha_\lambda = \lambda\tfrac{\beta}{4}e^{\frac{\beta}{4}u} du + \Re[(e^{-i\alpha_\lambda(u)}-1)(dZ-i\delta du)],\qquad \lambda\in \R, t\in(-\infty,\infty) 
\end{align}

\begin{align}\label{eq:SDE_initial}
    \lim_{u\to-\infty} \alpha_\lambda(u)=0, \qquad \lambda\in \R. 
\end{align}

\begin{proposition}\label{prop:alphaSDE}Fix $\beta>0$ and $\Re \delta>-1/2$. 


\begin{enumerate}
    \item For any $\nu\in \R$ the system \eqref{eq:alphaSDE} has a unique strong solution $\alpha_{\lambda,\nu}(u)$ on $[\nu,\infty)$ with initial condition $\alpha_{\lambda,\nu}(\nu)=0$. For any given $\nu< u$, with probability one the function $\lambda\mapsto \alpha_{\lambda,\nu}(u)$ is analytic and strictly increasing with $\alpha_{0,\nu}(u)=0$ 

% \item Suppose that $\alpha$ and $\tilde \alpha$ are both strong solutions of \eqref{eq:alphaSDE} on $[u_0,\infty)$. Let $0\le \lambda_1\le \lambda_2$. Then if $P(\alpha_{\lambda_1}(u_0)\le \tilde \alpha_{\lambda_2}(u_0))=1$ then  $\alpha_{\lambda_1}(u)\le \tilde \alpha_{\lambda_2}(u)$ for all $u_0\le u$ with probability one. 
% \benedek{Maybe we don't need this?}

    \item There is a unique strong solution $\alpha_\lambda(u)$ of \eqref{eq:alphaSDE} on $(-\infty,\infty)$ satisfying the initial condition \eqref{eq:SDE_initial}, which  can be obtained as the a.s.~limit $\alpha_\lambda(u)=\lim\limits_{\nu\to -\infty} \alpha_{\lambda,\nu}(u)$. The limit is monotone increasing for $\lambda>0$ and monotone decreasing for $\lambda<0$. 
    
    For any $u\in \R$ with probability one the function $\lambda\mapsto \alpha_\lambda(u)$ is analytic and strictly increasing in $\lambda$ with $\alpha_0(u)=0$. 

    Moreover, $\alpha_\lambda(u)$ satisfies the following scale-invariance property: for any fixed $s\in \R$ 
\begin{align}\label{eq:scale_inv}
    \left(\alpha_\lambda(t), t\in (-\infty,\infty)\right) \ed  \left(\alpha_{e^{\beta s/4}\lambda}(t-s), t\in (-\infty,\infty)\right). 
\end{align}  
\end{enumerate}
\end{proposition}

\begin{align}\notag
 d e^{i k\alpha_\lambda(u)}&=\frac{k(k+\delta)}{2} e^{i (k-1)\alpha_\lambda(u)} du+(i k \lambda \tfrac{\beta}{4}e^{\frac{\beta}{4}u}-k^2) e^{i k \alpha_\lambda(u)} du+\frac{k(k-\delta)}{2} e^{i (k+1)\alpha_\lambda(u)}du  \\
 &\qquad -\frac{1}{2} i  k (e^{i \alpha_\lambda(u)}-1) e^{i(k-1)\alpha_\lambda(u)} dZ+\frac{1}{2} i  k (e^{i \alpha_\lambda(u)}-1) e^{i k \alpha_\lambda(u)} d\bar Z. \label{eq:QkSDE}
\end{align}

\begin{align}\label{eq:q_coeff_1}
    \mathbf{s}_0&=-\frac{n+1}{2}\mathbf{A}_n^{-1}\mathbf{e}_n, \qquad \mathbf{s}_k=i \left(k \mathbf{I}-\tfrac{4}{\beta} \mathbf{A}_n\right)^{-1} \mathbf{B}_n \mathbf{s}_{k-1}, \qquad k\ge 1.
\end{align}

\begin{lemma}\label{lem:coeff_s} Consider the functions $ \mathbf{s}_{j,\nu}(u), j\ge 0, u\ge \nu$ defined in \eqref{eq:s_nu_def1}. For every $\nu\in \R$, these are well-defined, and they satisfy the  ODEs given in \eqref{eq:coeff_s_ode1} and   \eqref{eq:coeff_s_ode2}. Moreover, we have the following norm bounds for $\nu\le u$:
\begin{align}\label{eq:coeff_norm_bound}
        \|\mathbf{s}_{k,\nu}(u)\|\le e^{\frac{\beta}{4} k u} \,  \kappa^k \prod_{j=1}^k\frac{n}{j+\frac{4}{\beta}n},
\end{align}
with some $\kappa>0$ (which only depends on $n$). 
Moreover, these functions satisfy 
\begin{align}\label{eq:coeff_nu_lim}
    \lim\limits_{\nu\to -\infty} \mathbf{s}_{k,\nu}(u)=e^{k \frac{\beta}{4}u} \mathbf{s}_{k}.
\end{align}
With $\mathbf{q}_{\nu}(u)$ given by \eqref{eq:q_nu_series} we also have \eqref{eq:q_nu_lim}.
\end{lemma}

\begin{theorem}\label{thm:even_beta}
 The vector-valued function $\mathbf{q}(\lambda)$ satisfies the ODE
 \begin{align}\label{eq:ode_q_vec}
     \frac{\beta}{4} \lambda \mathbf{q}'(\lambda) = (i \frac{\beta}{4}\lambda  \mathbf{B}_n   +\mathbf{A}_n) \mathbf{q}(\lambda)+\frac{n+1}{2} \mathbf{e}_n, \qquad \mathbf{q}(0)=\mathbf{f}_n.
\end{align}
The functions $\lambda\mapsto q_k(\lambda)$ can be extended to $\CC$ as entire functions  so that \eqref{eq:ode_q_vec} holds for all $\lambda\in \CC$. With this extension, we have for all $\lambda\in \CC$
\begin{align}\label{eq:q_expansion}
    \mathbf{q}(\lambda)=\sum_{j=0}^\infty \mathbf{s}_j \lambda^j.
\end{align} 
\end{theorem}

\begin{proof}[Proof of Theorem \ref{thm:even_beta}]
For a fixed $\nu<0$ consider the unique strong solution $\alpha_{\lambda,\nu}(u)$ of \eqref{eq:alphaSDE} on $[\nu,\infty)$ with initial condition $\alpha_{\lambda,\nu}(\nu)=0$ and $\delta=n$. Set 
\begin{align}
    q_{k,\nu}(u,\lambda)=E[e^{i k \alpha_{\lambda,\nu}(u)}].
\end{align}
Fix $\lambda\in \R$. By Proposition \ref{prop:alphaSDE} we have  $\lim\limits_{\nu\to -\infty} \alpha_{\lambda,\nu}(u)=\alpha_\lambda(u)$ almost surely, hence by the bounded convergence theorem we have 
\begin{align}
    \lim_{\nu\to -\infty}  q_{k,\nu}(0,\lambda)=q_k(\lambda). 
\end{align}
The process $e^{i k \alpha_{\lambda,\nu}(u)}$ satisfies the same SDE as $e^{i k \alpha_\lambda(u)}$, see \eqref{eq:QkSDE}. For a fixed $\lambda$, the coefficients of this SDE are bounded on $[\nu,0]$, hence the Brownian terms do not contribute when we take expectations in the integrated form of the equation. 
% It\^o's formula gives  
% \begin{align}
%         \label{eq:Q_sde}
%     dQ_\nu(u,\lambda)=&\frac{\delta+1}{2}du+(i \lambda f(u)-1) Q_\nu du+\frac{1}{2}(1-\delta) Q_\nu^2 du\\
%     &\notag \qquad-\frac{1}{2} (Q_\nu-1) (i db_1 Q_\nu-i db_1+db_2 Q_\nu+db_2)\\
%     dQ_\nu^k(u,\lambda)=&\frac{k}{2} (\delta +k)Q_\nu^{k-1}+(-k^2 + i k\lambda f_\beta(u))Q_\nu^{k}+\frac{1}{2} k (k-\delta )Q_\nu^{k+1}du\nonumber\\
%    &+\frac{1}{2} i k Q_\nu^{k-1} (1-Q_\nu)^2 db_1+\frac{1}{2} k Q_\nu^{k-1} \left(Q_\nu^2-1\right)db_2
% \end{align}
% The coefficients of these SDEs are bounded for $\nu\le u\le t$ (for any fixed $t>\nu$), hence we can take expectations to get 
This leads to the following ODE system for $u\in [\nu,\infty)$.
\begin{align}\label{eq:ODEq1}
    \partial_u q_{1,\nu}&=\frac{n+1}{2}+(i \lambda \tfrac{\beta}{4}e^{\frac{\beta}{4}u}-1) q_{1,\nu}+\frac{1}{2}(1-n)q_{2,\nu},\\\label{eq:ODEq2}
    \partial_u q_{k,\nu}&=\frac{k}{2} (n +k) q_{k-1,\nu}+(-k^2 + i k \lambda \tfrac{\beta}{4}e^{\frac{\beta}{4}u})q_{k,\nu}+\frac{1}{2} k (k-n )q_{k+1,\nu}, \qquad 2\le k\le n,\\[3pt]
   &q_{k,\nu}(\nu,\lambda)=1, \qquad 1\le k\le n. 
\end{align}
We did not denote the dependence on $u,\lambda$ in \eqref{eq:ODEq1} and \eqref{eq:ODEq2}. 
Since the coefficient 
of $q_{n+1,\nu}$ in \eqref{eq:ODEq2} is 0,   we obtained a closed ODE system.
Using the notation
\[
\mathbf{q}_{\nu}(u,\lambda)=[q_{1,\nu}(u,\lambda), \dots, q_{n,\nu}(u,\lambda)]^T,
\]
we can write it as 
\begin{align}\label{eq:q_nu_ODE}
     \partial_u \mathbf{q}_\nu(u,\lambda) &= (i \tfrac{\beta}{4}e^{\frac{\beta}{4}u} \mathbf{B}_n \lambda  +\mathbf{A}_n) \mathbf{q}_\nu(u,\lambda) +\frac{n+1}{2}\mathbf{e}_n, \qquad
     \mathbf{q}_{\nu}(\nu,\lambda) = \mathbf{f}_n.
\end{align}
For a fixed $\lambda\in \CC$ the coefficients of this ODE system are bounded for $u\in [\nu,0]$, hence there is unique solution on $[\nu,0]$ for each $\lambda\in \CC$, and this solution is analytic in $\lambda$ for a fixed $u$. We claim that this unique solution is given by 
\begin{align}
    \mathbf{q}_\nu(u,\lambda)=\sum_{j=0}^\infty \mathbf{s}_{j,\nu}(u)\lambda^j,\label{eq:q_nu_series}
\end{align}
where the vector valued functions $\mathbf{s}_{j,\nu}(u),j\ge 0$ are  given by the recursion  
\begin{align}\label{eq:s_nu_def1}
    \mathbf{s}_{0,\nu}(u) &=\frac{n+1}{2}\int_\nu^u e^{(u-s)\mathbf{A}_n} \mathbf{e}_nds+e^{(u-\nu) \mathbf{A}_n }\cdot \mathbf{f}_n,\\
    \mathbf{s}_{k,\nu}(u) &= i \frac{\beta}{4}e^{\frac{\beta}{4} u}\int_\nu^ue^{n/2 (s-u)}e^{(u-s)\mathbf{A}_n}\mathbf{B}_n \mathbf{s}_{k-1,\nu}(s)ds.\label{eq:s_nu_def2}
\end{align}
Lemma \ref{lem:coeff_s} below shows that for $k\ge 0$ the $\CC^n$ valued function $\mathbf{s}_{k,\nu}(u)$ are bounded in norm by a $k$-dependent constant multiple of $e^{\frac{\beta}{4} k u}$, and that these functions satisfy the ODEs
\begin{align}\label{eq:coeff_s_ode1}
     \mathbf{s}_{0,\nu}'(u)&= \mathbf{A}_n  \mathbf{s}_{0,\nu}(u)+ \frac{n+1}{2}\mathbf{e}_n,\\
        \mathbf{s}_{k,\nu}'(u)&= \mathbf{A}_n \mathbf{s}_{k,\nu}(u)+ i \tfrac{\beta}{4}e^{\frac{\beta}{4}u}\mathbf{B}_n \mathbf{s}_{k-1,\nu}(u), \qquad k\ge 1.\label{eq:coeff_s_ode2}
\end{align}
From this it follows that the series given by  \eqref{eq:q_nu_series} can be differentiated in $u$ term-by-term and the resulting series is absolutely convergent. Moreover, the series in \eqref{eq:q_nu_series} satisfies the ODE system \eqref{eq:q_nu_ODE} and hence it is the unique solution of that system.  Lemma \ref{lem:coeff_s} also shows that for any $u>-\infty$ we have $\lim\limits_{\nu\to -\infty} \mathbf{s}_{k,\nu}(u)=e^{k \frac{\beta}{4}u} \mathbf{s}_{k}$ where $\mathbf{s}_{k}$ satisfies \eqref{eq:q_coeff_1}, and that 
\begin{align}\label{eq:q_nu_lim}
    \lim_{\nu\to -\infty} \mathbf{q}_{\nu}(u,\lambda)=\sum_{j=0}^\infty \mathbf{s}_j e^{\frac{\beta}{4} j u} \lambda^j.
\end{align}
From this, we obtain \eqref{eq:q_expansion}, and direct differentiation (which is allowed by Lemma \ref{lem:coeff_s}) gives \eqref{eq:ode_q_vec} as well.
\end{proof}   

\end{document}
